\documentclass[12pt,reqno]{amsart}
\DeclareSymbolFont{AMSb}{U}{msb}{m}{n}
\DeclareMathSymbol{\N}{\mathbin}{AMSb}{"4E}
\DeclareMathSymbol{\Z}{\mathbin}{AMSb}{"5A}
%\DeclareMathSymbol{\R}{\mathbin}{AMSb}{"52}
%\DeclareMathSymbol{\Q}{\mathbin}{AMSb}{"51}
%\DeclareMathSymbol{\I}{\mathbin}{AMSb}{"49}
\DeclareMathSymbol{\C}{\mathbin}{AMSb}{"43}

\def\sech{{\mathrm{sech}\,}}

\font\fff= eufm10 scaled \magstep 1
\def\?#1?{\hbox{\fff #1}}

\parindent=0pt

\def\pidagger{\pi^\ddag}
\def\pistar{\pi^\ast}
\def\pitilde{\tilde\pi}
\def\pihat{\hat\pi}
\def\pihatstar{\hat\pi^\ast}
\def\picheck{\check\pi}
\def\xid{\xi^\ddag}
\def\xis{\xi^\ast}
\def\xicheck{\check\xi}
\def\ad #1.{{{\rm ad}\,#1}}
\def\X{{\hat X}}
\def\Y{{\hat Y}}
\def\D{{\hat D}}
\def\R{{\cal R}}
\def\V{{\hat V}}
\def\W{{\hat W}}
\def\y{{\hat y}}
%\def\C{{\bf C}}
\def\e{{\bf e}}
\def\1{{\bf 1}}
\def\tfrac#1#2{\frac{\textstyle #1}{\textstyle #2}}
\def\eqref#1{(\ref{#1})}

\begingroup 
\newtheorem{bigthm}{Theorem}                 % Numbered separately, as A,B, etc.
\newtheorem{thm}{Theorem}[section]           % Numbered within each section
\newtheorem{lemma}[thm]{Lemma}
\newtheorem{proposition}[thm]{Proposition}
\newtheorem{theorem}[thm]{Theorem}
\newtheorem{corollary}[thm]{Corollary}
\newtheorem{definition}[thm]{Definition}
\newtheorem{notation}{Notation.}
\newtheorem{terminology}{Terminology}
\newtheorem{remark}[thm]{Remark. }             % Numbered along with thm
\newtheorem{example}[thm]{Example}           % Numbered along with thm
\endgroup

\def\qed{\qquad\framebox[7pt] \medskip \noindent}
\def\pqed{\rightline{\qed} \medskip \noindent}

\def\th{\thinspace}

\def\half{{\textstyle\frac{1}{2}}}
\def\avg#1.{\langle#1\rangle}
\def\vstrut{\phantom{\biggm|}}

\renewcommand{\thebigthm}{\Alph{bigthm}}  % Number as "Theorem A."
\renewcommand{\thenotation}{}  
\renewcommand{\theterminology}{}  
\renewenvironment{remark}{\bigskip \rm\ {\bf Remark. }}{\bigskip}

\renewenvironment{abstract}{\begin{quote}{\bf Abstract.\
}\small}{\end{quote}\bigskip}

\def\vstrut{\phantom{\biggm|}}

\title{Estimating the sum of the absolute values of the entries in the Krawtchouk matrices}

\author{Philip Feinsilver}\address{Department of Mathematics
Southern Illinois University 
 Carbondale, IL. 62901, U.S.A.}

\date{}

\begin{document} 
\maketitle
\thispagestyle{empty}
Start with the symmetry of $S=\Phi B$, so $\Phi B=B\Phi^*$. With $\Phi^2=2^NI$, we have
$$\Phi B\Phi^*=\Phi^2 B=2^NB$$
Rescaling by $2^N$ and adjusting factors of $\sqrt{B}$, we get
$$2^{-N/2}\sqrt{B}^{-1}\Phi\sqrt{B} 2^{-N/2}\sqrt{B}\Phi^*\sqrt{B}^{-1}=I$$
i.e., $2^{-N/2}\sqrt{B}^{-1}\Phi\sqrt{B}$ is orthogonal. Taking squared norms of the rows yields
$$\sum_j \Phi_{ij}^2\binom{N}{j}=2^N\binom{N}{i}$$
Similary, the columns yield
$$\sum_i \Phi_{ij}^2\binom{N}{i}^{-1}=2^N\binom{N}{j}^{-1}$$

Schwartz' inequality for the sums over $j$ yield the estimate
$$\sum_j |\Phi_{ij}|\,{\textstyle\sqrt{\binom{N}{j}}/\sqrt{\binom{N}{j}}}\le \sqrt{\sum_j \textstyle\binom{N}{j}^{-1}}\,
\sqrt{\textstyle\binom{N}{i}2^N}$$
while the sums over $i$ yield
$$\sum_i |\Phi_{ij}|\,{\textstyle\sqrt{\binom{N}{i}}/\sqrt{\binom{N}{i}}}\le \sqrt{2^N}\,\sqrt{2^N/\textstyle\binom{N}{j}}
=2^N\binom{N}{j}^{-1/2}$$

Since the column sums for $j=0$ and $j=N$ are $2^N$, 
$$2^{N+1}\le\sum_{i,j} |\Phi_{ij}|\le 2^N\binom{N}{j}^{-1/2}=2^N\bigl(2+\sum_{0<j<N}\binom{N}{j}^{-1/2}\bigr)$$
For $N>9$, say, we have, by symmetry and unimodality of the binomial coefficients,
\begin{align*}
\sum_{0<j<N}\binom{N}{j}^{-1/2}&\le 2\left[\frac{1}{\sqrt{N}}+\frac{\sqrt{2}}{\sqrt{N(N-1)}}+\frac{\sqrt{6}\,(N/2-1)}{\sqrt{N(N-1)(N-2)}}\right]\cr
&=2\left[\frac{1}{\sqrt{N}}+\frac{\sqrt{2}}{\sqrt{N(N-1)}}+\frac{\sqrt{6(N-2)}/2}{\sqrt{N(N-1)}}\right]\cr
&\le\frac{2+1+\sqrt{6}/2}{\sqrt{N}}\le\frac{5}{\sqrt{N}}
\end{align*}
using $2\sqrt{2}<\sqrt{N-1}$. This gives
$$
2^{N+1}\le\sum_{i,j} |\Phi_{ij}|\le 2^N\,\bigl(2+\frac{5}{\sqrt{N}}\bigr)$$
or
$$\sum_{i,j} |\Phi_{ij}|\sim 2^{N+1}\,, \qquad\text{as } N\to\infty
$$
The second column sums to $\displaystyle \binom{N}{N/2}$ for even $N$, and generally, to
$\displaystyle \binom{N-1}{\lfloor(N-1)/2\rfloor}$. This improves the lower bound to
$$2^{N+1}+2\,\binom{N-1}{\lfloor(N-1)/2\rfloor}$$
Stirling estimates give
$$2^{N+1}\,\sqrt{\frac{2}{\pi N}}$$
for the additional term.





\end{document}












